\section{Clifford Algebra}
Let us now start with the definition of the Clifford Algebra. Well, there are different ways of definining the Clifford Algebra, through a pure algebraic approach, through the exterior algebra, through tensor algebra, etc. All of these are mostly equivalent, however the definition of the Clifford Algebra as a quotient algebra of the tensor algebra generalises the notion to rings also, in addition to field. \\[0.2cm]
I understand that all the hoity-toity things that I said above does not make any sense. Thus, we will try to forsake the rigour and understand Clifford algebra as intuitively as possible. The foremost thing to define a Clifford algebra is a vector space $V$ over a field $K$, equipped with a quadratic form $Q:V\to K$.\\[0.2cm] For now, think of quadratic forms as a polynomial expression with terms of degree 2 (eg. $x^2,y^2,xy\ldots$). What $Q$ does it, it takes some $\vb{v}\in V$ and then represents $\vb{v}$ in some basis $\SB$ of $V$ and then puts the coordinates of $\vb{v}$ in that basis into the polynomial expression to generate some element in the field $K$. \\[0.2cm]
Here we will consider Clifford algebras as an extension of the \textit{exterior algebra} which was introduced by Grassman.
\subsection{Exterior Algebra}
This is the algebra introduced by \textit{Grassmann} where the bilinear binary operation of the algebra is the \textit{wedge product} or the \textit{exterior product}. The exterior product of vectors comes from differential geometry and a basic introduction can be found \href{https://raw.githubusercontent.com/LoneWolf1304/Notes/main/Tensor%20Analysis/main.pdf}{here}!